\begin{tikzpicture}[figure picture,x=10mm,y=10mm,font=\figurefont\fontsize{7.8}{9.8}\selectfont]
\path[use as bounding box] (0,0) rectangle (16.4,4.7);
% (a) 弱收敛但 TV=1
\node[figure heading,anchor=west] at (0.4,4.35) {(a) $P_n=\delta_{1/n}\Rightarrow\delta_0$，但 $\operatorname{TV}(P_n,P)=1$};
\draw[FigureGrid,line width=.8pt,-{Stealth[length=1.6mm]}] (0.8,2.4)--(7.2,2.4);
\foreach \k/\xx in {1/3.0,2/2.1,3/1.65,5/1.35} \fill[FigureBlue] (\xx,2.4) circle[radius=.8mm];
\fill[FigureRed] (1.1,2.4) circle[radius=1.1mm];
\node[below,FigureMuted] at (1.1,2.22) {$0$};
\draw[FigureRed,line width=.7pt,{Stealth[length=1.4mm]}-{Stealth[length=1.4mm]}] (1.1,2.95)--(3.0,2.95);
\node[above,FigureMuted] at (2.05,3.0) {$1/n$};
\draw[FigureBlue,line width=.6pt,dashed,-{Stealth[length=1.4mm]}] (3.0,2.25) to[bend right=25] (1.25,2.32);
\node[figure note,align=left] at (3.9,1.35) {连续探针看到质量在位置上不断靠近；\\但事件 $\{0\}$ 概率始终为 $0$，全变差判为完全不同};
% 分隔
\draw[FigureGrid,line width=.5pt,dashed] (8.0,.3)--(8.0,4.2);
% (b) 远端小质量
\node[figure heading,anchor=west] at (8.3,4.35) {(b) $P_n=(1-n^{-p})\delta_0+n^{-p}\delta_n$：弱收敛，但 $W_p(P_n,P)=1$};
\draw[FigureGrid,line width=.8pt,-{Stealth[length=1.6mm]}] (8.7,2.4)--(15.9,2.4);
\fill[FigureRed] (9.1,2.4) circle[radius=1.4mm];
\node[below,FigureRed] at (9.1,2.2) {质量 $1-n^{-p}$};
\fill[FigureBlue] (14.9,2.4) circle[radius=.9mm];
\node[below,FigureBlue] at (14.9,2.2) {$n$ 处仅 $n^{-p}$};
\node[below,FigureMuted] at (9.1,2.62) {$0$};
\draw[FigureBlue,line width=.8pt,{Stealth[length=1.4mm]}-{Stealth[length=1.4mm]}] (9.1,3.15)--(14.9,3.15);
\node[above,FigureBlue] at (12.0,3.2) {距离 $n$，代价 $n^{-p}\cdot n^p=1$};
\node[figure note,align=left] at (12.2,1.25) {概率很小 $\times$ 距离很大 $=$ 常数代价；\\没有尾部的共同控制，Wasserstein 不收敛};
\end{tikzpicture}
